% Source: https://artofproblemsolving.com/community/c3h2346777_circles__geo
% Stable identity: aops-topic-2346777-q01
The point $(8,9)$ is on the line, since $24+36=60$. The radius to the point of tangency is perpendicular to the tangent line, and $(3,4)$ is perpendicular to the line $3x+4y=60$. So the centre is $(8+3t,\ 9+4t)$ for some number $t$, and the radius is its distance from $(8,9)$, namely $5|t|$.

A circle tangent to both axes has centre $(\pm r,\pm r)$, so we need $|8+3t|=|9+4t|=5|t|$.
\begin{itemize}
\item If $8+3t=9+4t$, then $t=-1$: the centre is $(5,5)$, and indeed $|5|=5=5|t|$, so $r=5$.
\item If $8+3t=-(9+4t)$, then $t=-\dfrac{17}{7}$: the centre is $\left(\dfrac57,-\dfrac57\right)$, at distance $\dfrac57$ from the axes but $5|t|=\dfrac{85}7$ from $(8,9)$, so this fails.
\end{itemize}

\begin{asy}
usepackage("amsmath");
size(6.5cm);
defaultpen(fontsize(10pt));
draw((-1,0)--(22,0),Arrow(5)); draw((0,-1)--(0,17),Arrow(5));
label("$x$",(22,0),S); label("$y$",(0,17),W);
draw((20,0)--(0,15),blue+linewidth(0.9));
draw(circle((5,5),5),red+linewidth(0.9));
draw((5,5)--(8,9),dashed); draw((5,5)--(5,0),dotted); draw((5,5)--(0,5),dotted);
pair u=unit((5,5)-(8,9)), v=unit((20,0)-(8,9)); draw((8,9)+0.8*u--(8,9)+0.8*u+0.8*v--(8,9)+0.8*v);
dot((5,5)); dot((8,9)); dot((5,0)); dot((0,5));
label("$(5,5)$",(5,5),SW,fontsize(8pt)); label("$(8,9)$",(8,9),NE,fontsize(8pt));
label("$3x+4y=60$",(15,3.75),NE,blue+fontsize(8pt));
label("$20$",(20,0),S,fontsize(8pt)); label("$15$",(0,15),W,fontsize(8pt));
\end{asy}

So the circle is
\[\boxed{(x-5)^2+(y-5)^2=25}.\]
Check: $(8-5)^2+(9-5)^2=25$, and the distance from $(5,5)$ to the line is $\dfrac{|15+20-60|}{5}=5$, the radius.
